\documentclass[12pt]{article}
\usepackage[utf8]{inputenc}
\usepackage[french]{babel}
\usepackage{amsmath, amssymb, amsthm}
\usepackage{geometry}
\geometry{a4paper, margin=1in}

\title{Analyse Rigoureuse de la Conjecture d'Erd\H{o}s-Straus}
\author{Charles EDOU NZE\thanks{Charles EDOU NZE, chercheur indépendant}}
\date{\today}

\begin{document}
\maketitle

\begin{abstract}
Cet article présente une analyse formelle de la conjecture d'Erd\H{o}s-Straus, en la décomposant en définitions axiomatiques fondamentales, recherche contextuelle de la littérature, et lemmes intermédiaires sans ellipse. Il fournit également une architecture stricte pour l'autoformalisation dans des assistants de preuve.
\end{abstract}

\section{Définitions Axiomatiques et Formulation}
Soit $\mathbb{N} = \{1, 2, 3, \ldots\}$ l'ensemble des entiers naturels non nuls.
La conjecture d'Erd\H{o}s-Straus énonce que pour tout entier $n \ge 2$, il existe des entiers strictement positifs $x, y, z \in \mathbb{N}$ tels que :
\begin{equation}
\frac{4}{n} = \frac{1}{x} + \frac{1}{y} + \frac{1}{z}
\end{equation}

Pour formaliser ceci, nous définissons l'ensemble des solutions pour un $n \ge 2$ donné :
\begin{equation}
S(n) = \left\{ (x, y, z) \in \mathbb{N}^3 \;\middle|\; \frac{4}{n} = \frac{1}{x} + \frac{1}{y} + \frac{1}{z} \right\}
\end{equation}
La conjecture est équivalente à l'énoncé : $\forall n \ge 2, S(n) \neq \emptyset$.

\section{Recherche de Littérature Contextuelle}
La conjecture d'Erd\H{o}s-Straus généralise le problème de la représentation en fractions égyptiennes. Il est connu que pour tout nombre rationnel positif $p/q$, il existe des représentations sous forme de sommes de fractions unitaires distinctes. Le cas spécifique de $4/n$ contraint le nombre de termes à trois.

Historiquement, Mordell (1967) a prouvé que la conjecture est vérifiée sauf potentiellement pour les nombres premiers $n$ congrus à $1, 11, 13, 17, 19, \text{ ou } 23 \pmod{24}$. De plus, Webb et d'autres ont établi des bornes en utilisant des méthodes de crible, montrant que le nombre d'exceptions jusqu'à $N$ est au plus de $O(N \exp(-c \log N / \log \log N))$, ce qui est étroitement lié à la distribution des nombres premiers dans les progressions arithmétiques.

Par analogie avec la résolution récente de la conjecture faible de Goldbach par Helfgott (2013), qui a employé la méthode du cercle pour gérer les combinaisons additives de nombres premiers, on pourrait considérer une approche analytique localisée similaire. Cependant, la nature diophantienne non linéaire de cette équation suggère que des méthodes géométriques sur des surfaces algébriques, semblables à celles utilisées par Manin (1986), pourraient être plus directement applicables pour borner l'ensemble des non-solutions.

\section{Lemmes et Stratégie de Preuve}

Nous décomposons le problème en classes modulaires. Une identité bien connue nous permet de traiter de nombreux cas.

\textbf{Lemme 1 (Identité pour des formes modulaires spécifiques) :} Si $n \equiv 2 \pmod{3}$, alors $S(n)$ est non vide.

\begin{proof}
Soit $n \ge 2$ et supposons $n \equiv 2 \pmod{3}$. Ceci implique qu'il existe un entier $k \ge 0$ tel que $n = 3k + 2$.
Nous substituons $n = 3k + 2$ dans l'expression cible $\frac{4}{n}$ et cherchons un développement.
Considérons les termes $x = n$, $y = k+1$, et $z = n(k+1)$. Nous devons vérifier que :
\begin{equation}
\frac{1}{n} + \frac{1}{k+1} + \frac{1}{n(k+1)} = \frac{4}{n}
\end{equation}
En partant du côté gauche, nous trouvons un dénominateur commun $n(k+1)$ :
\begin{equation}
\frac{1}{n} + \frac{1}{k+1} + \frac{1}{n(k+1)} = \frac{k+1}{n(k+1)} + \frac{n}{n(k+1)} + \frac{1}{n(k+1)}
\end{equation}
En sommant les numérateurs, nous obtenons :
\begin{equation}
\frac{k+1 + n + 1}{n(k+1)} = \frac{n + k + 2}{n(k+1)}
\end{equation}
La substitution de $n = 3k + 2$ dans le numérateur donne :
\begin{equation}
\frac{3k + 2 + k + 2}{n(k+1)} = \frac{4k + 4}{n(k+1)} = \frac{4(k+1)}{n(k+1)}
\end{equation}
Puisque $k \ge 0$, $k+1 \neq 0$, nous pouvons simplifier par $k+1$ au numérateur et au dénominateur :
\begin{equation}
\frac{4(k+1)}{n(k+1)} = \frac{4}{n}
\end{equation}
Ainsi, $(n, k+1, n(k+1)) \in S(n)$, ce qui prouve le lemme.
\end{proof}

\section{Architecture pour l'Autoformalisation}
Pour intégrer cette preuve dans un assistant comme Lean 4, la structure doit être explicitement typée :

\begin{verbatim}
import Mathlib.Data.Nat.Basic
import Mathlib.Data.Rat.Basic

-- Définitions Axiomatiques
def S (n : Nat) : Set (Nat * Nat * Nat) :=
  { p | 4 * (p.1 * p.2.1 * p.2.2) =
        n * (p.2.1 * p.2.2 + p.1 * p.2.2 + p.1 * p.2.1)
        /\ p.1 > 0 /\ p.2.1 > 0 /\ p.2.2 > 0 }

-- Type de la Conjecture
def erdos_straus_conjecture : Prop :=
  forall n : Nat, n >= 2 -> Not (S n = empty)

-- Type du Lemme 1
lemma lemma_mod3 (n : Nat) (h1 : n >= 2) (h2 : n % 3 = 2) : Not (S n = empty) :=
  sorry
\end{verbatim}

\end{document}
